\subsection{An exact all-weight reduction at 2}

Continue with the checked presentation and normal form at $2$.
Identify $H_{\mathrm{ab}}$ with
$\Lambda=\mathbb Z_2[[T]]$, where $\gamma=1+T$ and $[h]=1$.
For $i\geq0$ put
\[
 m_i=2i+3,\qquad r_i=5^{m_i}-1,\qquad
 f_i(T)=\prod_{j<i}(T-r_j),\qquad \delta_i=f_i(r_i).
\]
Thus $v_2(\delta_i)=3i+v_2(i!)$.

\begin{proposition}[K17-2-NEWTON-BASIS]\label{prop:newton-basis}
Suppose that for every $i\geq0$ the actual primitive odd
image ideal in weight $m_i$ attains the Newton lower bound
$v_2(\delta_i)$. Then there are primitive classes in those odd
weights that freely generate the entire integral tripod graded
Lie algebra over $\mathbb Z_2$.
This proposition isolates a sufficient lifting hypothesis. The finite
calculation above does not establish it; Theorem~\ref{thm:2-allweights}
below establishes it in every weight.
\end{proposition}

\begin{proof}
The linear weight-$m$ coordinate of an element with abelian
class $F(T)$ is $F(5^m-1)$ in the normalized rational
coordinates. Choose $g_i\in H$ of actual weight at least
$m_i$ attaining the indicated minimum. Its earlier odd linear
coordinates vanish. Successive division by the monic factors
$T-r_j$ therefore gives
\[
 [g_i]=f_i(T)b_i(T).
\]
Its weight-$m_i$ coordinate is $\delta_i b_i(r_i)$, so
$b_i(r_i)$ is a unit. Since $r_i\in4\mathbb Z_2$, this is
equivalent to $b_i(0)$ being a unit, hence
$b_i\in\Lambda^\times$. Conversely, any lift of such an
abelian class having actual weight at least $m_i$ attains the
minimum. This proves the precise equivalence of the two
lifting formulations.

The classes $f_i b_i$ form a compact $\mathbb Z_2$-basis of
$\Lambda$. To check this, consider
\[
 B:\prod_{i\geq0}\mathbb Z_2\longrightarrow\Lambda,
       \qquad (a_i)\longmapsto\sum_{i\geq0}a_i f_i b_i.
\]
The series defines a continuous map. Indeed $r_j\in4\mathbb Z_2$
implies that the coefficient of $T^k$ in $f_i$ is divisible by
$2^{2(i-k)}$ for $i>k$; multiplication by an integral $b_i$
retains the required coefficientwise convergence. Equivalently,
the columns have increasing $(2,T)$-adic order. Modulo $2$,
the columns are $T^i\overline b_i(T)$, with unit leading
coefficient. Successive coefficient comparison makes this
triangular formal-series map an isomorphism. Lift modulo $2$,
subtract, divide by $2$, and repeat to prove surjectivity of
$B$. These corrections converge in the product topology.
The same argument applied to a kernel element shows that all
its coefficients are divisible by every power of $2$ and hence
zero. Thus $B$ is an isomorphism of compact modules.

The $g_i$ converge to $1$ in $H$. Their images have weights
tending to infinity, so they converge to $1$ in every fixed
rational weight truncation. The checked faithful rational
realization embeds compact $H$ into a Hausdorff inverse limit
and is a homeomorphism onto its image. The convergence therefore
holds in $H$ itself.

Take the free pro-$2$ group on the pointed compact sequence
$\{g_i:i\geq0\}\cup\{1\}$, where $g_i\to1$. It is the inverse
limit of the finite-rank free groups with successive tail
generators killed: every map of the convergent sequence into a
finite group kills its tail. This verifies the universal
property and also gives Magnus injectivity by finite-rank
projection. Its map to $H$
is an isomorphism on abelianizations by the preceding compact
basis calculation. Hence it is Frattini-surjective and thus
surjective. Projectivity of the free pro-$2$ group $H$ supplies
a section. The section has closed Frattini-surjective image in
the source, so its image is the whole source. Consequently the
$g_i$ form a compact free basis. This uses projectivity and
Frattini theory, not infinite-rank Hopficity.

Give the integral Magnus variable $X_i=g_i-1$ weight $m_i$.
A nonidentity group element has a first nonzero weighted Magnus
term. That term is primitive for the associated graded tensor
coproduct, because the group element is grouplike. The primitive
submodule of the integral tensor algebra is the free
$\mathbb Z_2$ Lie algebra: the rational primitive statement and
the unit-triangular Hall--Lyndon basis show that this Lie module
is saturated. Each arithmetic weight involves only finitely
many variables and words.

Let $\sigma_i$ be the actual weight-$m_i$ leading class of
$g_i$. Its $s_{m_i}$ indecomposable coefficient is
$\delta_i b_i(r_i)\ne0$. Every other term is a bracket in
strictly lower odd weights. Thus the homogeneous substitution
$X_i\mapsto\sigma_i$ is a triangular automorphism of the
rational free Lie algebra. It cannot kill a first nonzero
integral Magnus term. Therefore actual rational weight equals
abstract weighted Magnus degree, and every integral leading
class is a $\mathbb Z_2$ Lie polynomial in the $\sigma_i$.
Conversely every Hall bracket is realized by a group
commutator. This proves integral freeness for the rational
pullback filtration. Proposition~\ref{prop:outer-sat} identifies
it with the actual integral outer filtration.
\end{proof}

The all-weight problem is therefore reduced to the existence
of these lifts, or equivalently attainment of the Newton
minimum in every odd weight. This gate is closed by
Theorem~\ref{thm:2-allweights} below. A raw Newton word can have lower
even or decomposable terms, whose removal requires actual
integral commutators. The scalar abelian calculation does not
perform those corrections. Nor does a finite list of attained
odd minima certify all intervening lattices: later Newton
directions may still project to lower commutator weights.

Complex conjugation does not make every $\gamma$-conjugate of
$h$ exactly inverted, since $c\gamma c=h\gamma$. In particular,
the rational diagonal action $\tau$ differs from actual $D$ by
the half-power conjugation $\operatorname{Ad}(\exp(a/2))$.
Its integrality cannot be assumed. The finite certificate uses
actual $D$ and retains every denominator.
